\section{Fourier estimates, product projections, and cross blocks}
\label{sec:fourier}

Throughout this section the coordinates of \(X=(X_1,\ldots,X_n)\) are
independent, and their probability densities satisfy \(0\leq f_i\leq1\).
We use the Fourier transform
\[
 \widehat h(\xi)=\int_{\mathbb R^d}e^{i\langle \xi,x\rangle}h(x)\,dx,
 \qquad
 h(x)=(2\pi)^{-d}\int_{\mathbb R^d}
                 e^{-i\langle \xi,x\rangle}\widehat h(\xi)\,d\xi.
\]
Statements about essential suprema do not depend on the choice of a
density representative.  When the characteristic function is integrable,
Fourier inversion supplies a continuous density vanishing at infinity.

\subsection{Classical analytic inputs and an inversion formula}

We use the geometric rank-one Brascamp--Lieb inequality
\citep{BL1976,LPP2016}: if \(v_j\) are unit vectors in \(\mathbb R^d\),
\(c_j>0\), and \(\sum_jc_jv_jv_j^{\mathsf T}=I_d\), then
\begin{equation}\label{four:geometric-bl}
 \int_{\mathbb R^d}\prod_jH_j(\langle v_j,z\rangle)^{c_j}\,dz
 \leq\prod_j\left(\int_{\mathbb R}H_j\right)^{c_j}.
\end{equation}
The functions \(H_j\) are arbitrary nonnegative measurable functions.
We also use Lieb's Gaussian exhaustion theorem \citep{Lieb1990}.
For nonzero row vectors \(b_j\), positive exponents \(c_j\) with
\(\sum_jc_j=d\), the best constant in the corresponding rank-one
Brascamp--Lieb inequality is
\begin{equation}\label{four:gaussian-bl}
 \sup_{\lambda_j>0}
 \left[
 \frac{\prod_j\lambda_j^{c_j}}
      {\det\left(\sum_jc_j\lambda_jb_jb_j^{\mathsf T}\right)}
 \right]^{1/2}.
\end{equation}
The supremum may be infinite in general.  The determinant estimates
below show finiteness for the data considered here.

The sharp Hausdorff--Young and convolution inequalities
\citep{Beckner1975} are used with the following normalization:
\begin{align}
 \|\widehat h\|_{p'}&\leq
 (2\pi)^{1/p'}C_p\|h\|_p,
 &C_p&=p^{1/(2p)}(p')^{-1/(2p')},\quad 1\leq p\leq2,
 \label{four:hy}\\
 \|h_1*h_2\|_r&\leq\frac{C_pC_q}{C_r}
                     \|h_1\|_p\|h_2\|_q,
 &&\frac1p+\frac1q=1+\frac1r .
 \label{four:young}
\end{align}
In the convolution inequality \(1\leq p,q,r\leq\infty\), and
each \(C_t\) is given by the same conjugate-exponent formula.
Endpoint constants have their limiting values.  In particular
\eqref{four:hy} applies to complex functions.

\begin{lemma}[Integrable power envelopes]\label{four:inversion}
Let \(Y\) be real and suppose, for \(a>0\) and \(b>1\), that
\[
 |\varphi_Y(t)|\leq\min\{1,(a/|t|)^b\}.
\]
Then \(Y\) has a continuous density vanishing at infinity and
\begin{equation}\label{four:inversion-bound}
 \|p_Y\|_\infty\leq\frac{ab}{\pi(b-1)}.
\end{equation}
More generally, a bound \(K|t|^{-b}\) can be used with \(a=K^{1/b}\).
\end{lemma}
\begin{proof}
Splitting the integral at \(a\) gives
\[
 \int_{\mathbb R}|\varphi_Y(t)|\,dt
 \leq2a+2a^b\int_a^\infty t^{-b}\,dt
 =\frac{2ab}{b-1}.
\]
The inverse Fourier transform of \(\varphi_Y\) is continuous and vanishes
at infinity, and is the density of its probability law.  Its supremum
is at most \((2\pi)^{-1}\|\varphi_Y\|_1\).
\end{proof}

\subsection{Densities of projections and modulated scalar transforms}

\begin{lemma}[Product projection densities]\label{four:projection}
Let \(U\) be an \(n\times k\) real matrix with
\(U^{\mathsf T}U=I_k\), and put
\(c_i=\|U^{\mathsf T}e_i\|^2\).  Then \(U^{\mathsf T}X\) has a
density satisfying the explicit estimate
\begin{equation}\label{four:projection-elementary}
 \|p_{U^{\mathsf T}X}\|_\infty
 \leq\prod_{i:c_i<1}(1-c_i)^{-(1-c_i)/2}
 \leq e^{k/2}.
\end{equation}
The sharper classical product-marginal theorem \citep{LPP2016} gives
\begin{equation}\label{four:projection-sharp}
 \|p_{U^{\mathsf T}X}\|_\infty\leq2^{k/2}.
\end{equation}
For any rank-\(k\) map \(T:\mathbb R^n\to\mathbb R^k\), consequently,
\begin{equation}\label{four:linear-density}
 \|p_{TX}\|_\infty
 \leq\frac{2^{k/2}}{\sqrt{\det(TT^{\mathsf T})}}
 \leq\frac{e^{k/2}}{\sqrt{\det(TT^{\mathsf T})}}.
\end{equation}
\end{lemma}
\begin{proof}
To prove \eqref{four:projection-elementary}, complete \(U\)
to an orthogonal matrix \((U,W)\), with \(W\) having \(n-k\) columns.
Writing \(w_i^{\mathsf T}\) for its rows and \(b_i=\|w_i\|^2=1-c_i\),
Fubini's theorem gives a marginal density
\[
 p(y)=\int_{\mathbb R^{n-k}}
             \prod_i f_i((Uy)_i+\langle w_i,z\rangle)\,dz.
\]
Terms with \(b_i=0\) are at most one and can be omitted.  The other
vectors \(v_i=w_i/\sqrt{b_i}\) satisfy
\(\sum_i b_iv_iv_i^{\mathsf T}=I_{n-k}\).  Apply
\eqref{four:geometric-bl} to
\[
 H_i(s)=f_i((Uy)_i+\sqrt{b_i}s)^{1/b_i}.
\]
Because \(b_i\leq1\), \(f_i^{1/b_i}\leq f_i\), so
\[
 \left(\int H_i\right)^{b_i}
 \leq b_i^{-b_i/2}.
\]
This proves the product bound uniformly in \(y\).  The entropy
inequality
\begin{equation}\label{four:entropy}
 -(1-x)\log(1-x)\leq x,\qquad 0\leq x\leq1,
\end{equation}
with the endpoint interpreted continuously, shows that this product
is at most \(e^{\sum_i c_i/2}=e^{k/2}\).
If \(k=n\), the density is simply the orthogonal image of the product
density and is at most one.

The product-marginal theorem of \citet[Theorem 1.1]{LPP2016}
gives, for \(k<n\), the bound
\[
 \min\left\{
       \left(\frac{n}{n-k}\right)^{(n-k)/2},\,2^{k/2}
      \right\}.
\]
Its hypotheses hold for the independent product densities with cap
one considered here; see also \citet{RV2015} for a dimension-free
projection theorem.
Finally, write \(T=(TT^{\mathsf T})^{1/2}U^{\mathsf T}\) and apply
the change-of-variables formula to the invertible first factor.
\end{proof}

\begin{lemma}[Modulated scalar Fourier estimate]\label{four:lp}
Let \(f\) be a probability density with \(0\leq f\leq1\), and let
\(\omega\) be measurable with \(|\omega|\leq1\).
For \(h=f\omega\) and every real \(p\geq2\),
\begin{equation}\label{four:lp-sharp}
 \int_{\mathbb R}|\widehat h(\xi)|^p\,d\xi
 \leq \frac{2\pi}{\sqrt p}
       \left(\frac{p}{p-1}\right)^{(p-1)/2}
 \leq\frac{2\pi\sqrt e}{\sqrt p}
 \leq\frac{4\pi}{\sqrt p}.
\end{equation}
In particular the estimate is uniform in every measurable phase
\(\omega(x)=e^{iV(x)}\), and in translations of the Fourier variable.
\end{lemma}
\begin{proof}
Set \(p'=p/(p-1)\).  Since \(|h|\leq f\leq1\) and
\(\int f=1\), we have \(\|h\|_{p'}\leq1\).  Apply
\eqref{four:hy} with input exponent \(p'\), then raise to the
power \(p\):
\[
 \|\widehat h\|_p^p
 \leq2\pi(p')^{p/(2p')}p^{-1/2}
 =2\pi p^{-1/2}
       \left(1+\frac1{p-1}\right)^{(p-1)/2}.
\]
The inequality \(\log(1+x)\leq x\) proves the second bound, and
\(\sqrt e<2\) proves the last one.  Translation invariance of Lebesgue
measure proves the assertion about Fourier shifts.
\end{proof}

\begin{lemma}[Difference densities and Fourier powers]\label{four:scalar-norms}
Let \(f\) be as in Lemma~\ref{four:lp}, put
\(\widetilde f(x)=f(-x)\) and \(g=f*\widetilde f\).  Then
\(g\) is a probability density bounded by one, and, for \(0<c\leq1\),
\begin{equation}\label{four:difference-norm}
 \|g\|_{1/c}\leq
 2^c c^{c/2}(1+c)^{-(1+c)/2}.
\end{equation}
If \(q\geq1\) and \(0<d\leq q/2\), then
\begin{equation}\label{four:fourier-power-norm}
 \big\||\widehat f|^q\big\|_{1/d}
 \leq (2\pi)^d q^{q/2-d}(q-d)^{-(q-d)/2}d^{d/2}.
\end{equation}
\end{lemma}
\begin{proof}
The identity defining \(g\) gives \(\int g=1\) and
\(g(x)\leq\|f\|_\infty\int\widetilde f\leq1\).
Use \eqref{four:young} with both input exponents \(2/(1+c)\)
and output \(1/c\).  Their norms are at most one.  Substitution in
\(C_{2/(1+c)}^2/C_{1/c}\) gives
\(2^cc^{c/2}(1+c)^{-(1+c)/2}\), proving
\eqref{four:difference-norm}.  This remains valid at \(c=1\) by
continuity.
For the second estimate, the Fourier exponent \(q/d\) is at least
two and its conjugate input exponent is \(q/(q-d)\).  Apply
\eqref{four:hy} and raise to the power \(q\).  Since the input
norm is at most one, the resulting constant is
\[
 (2\pi)^d
 \left(\frac q{q-d}\right)^{(q-d)/2}
 \left(\frac qd\right)^{-d/2},
\]
which is \eqref{four:fourier-power-norm}.
\end{proof}

\subsection{Fresnel linearization with separable phases}

\begin{lemma}[Regularized Fresnel formula]\label{four:fresnel}
Let \(\Lambda\) be an invertible symmetric \(r\times r\) real matrix,
\(t\ne0\), and \(\varepsilon>0\).  Then
\begin{align}
 K_{\varepsilon,t}(\xi)
 &:=\int_{\mathbb R^r}
  e^{-\varepsilon|y|^2+it y^{\mathsf T}\Lambda y-i\langle\xi,y\rangle}\,dy
 \nonumber\\
 &=\pi^{r/2}\det(\varepsilon I-it\Lambda)^{-1/2}
       \exp\left[-\frac14
             \xi^{\mathsf T}(\varepsilon I-it\Lambda)^{-1}\xi\right],
 \label{four:fresnel-formula}
\end{align}
where the square root is the product of the branches having positive
real part in the scalar Gaussian factors.  Moreover,
\begin{equation}\label{four:fresnel-kernel}
 |K_{\varepsilon,t}(\xi)|
 \leq\pi^{r/2}|t|^{-r/2}|\det\Lambda|^{-1/2}.
\end{equation}
If \(g,\widehat g\in L^1(\mathbb R^r)\), it follows that
\begin{equation}\label{four:fresnel-integral}
 \left|\int_{\mathbb R^r}e^{it y^{\mathsf T}\Lambda y}g(y)\,dy\right|
 \leq(2\pi)^{-r}\|\widehat g\|_1
       \pi^{r/2}|t|^{-r/2}|\det\Lambda|^{-1/2}.
\end{equation}
\end{lemma}
\begin{proof}
Diagonalize \(\Lambda\) orthogonally.  Each one-dimensional integral
is
\[
 \int_{\mathbb R}e^{-(\varepsilon-it\lambda)y^2-i\eta y}\,dy
 =\sqrt\pi(\varepsilon-it\lambda)^{-1/2}
          e^{-\eta^2/[4(\varepsilon-it\lambda)]}.
\]
The Gaussian identity for a positive real coefficient extends to the
right half-plane by analytic continuation, with domination on compact
subsets; this proves the displayed identity for every \(\varepsilon>0\).
Multiplying the scalar factors proves \eqref{four:fresnel-formula}.
The real part of \((\varepsilon-it\lambda)^{-1}\) is
\(\varepsilon/(\varepsilon^2+t^2\lambda^2)>0\), so the exponential
factors have modulus at most one.  Also
\(|\varepsilon-it\lambda|^{-1/2}\leq|t\lambda|^{-1/2}\).
These observations prove \eqref{four:fresnel-kernel}.

Insert the inverse Fourier transform of \(g\) in the integral with
the Gaussian cutoff.  Fubini is justified by
\(\widehat g\in L^1\) and the integrable Gaussian.  Bound its inner
integral by \eqref{four:fresnel-kernel}.  Finally let
\(\varepsilon\downarrow0\); the left side converges by domination
with \(|g|\).  This proves \eqref{four:fresnel-integral}.
\end{proof}

\begin{theorem}[Diffuse quadratic components and separable phases]
\label{four:diffuse}
Let \(U^{\mathsf T}U=I_r\), let
\(c_i=\|U^{\mathsf T}e_i\|^2\leq1/2\), and let
\(\Lambda\) be invertible and symmetric.  For arbitrary real
measurable functions \(V_i\), finite almost surely, define
\[
 Y=(U^{\mathsf T}X)^{\mathsf T}\Lambda(U^{\mathsf T}X)
          +\sum_iV_i(X_i),\qquad
 B(U)=\prod_{i:c_i>0}(1-c_i)^{-(1-c_i)/2}.
\]
Writing \(\Delta=|\det\Lambda|\), for every \(r\geq1\),
\begin{equation}\label{four:diffuse-decay}
 |\varphi_Y(t)|
 \leq\min\{1,B(U)\pi^{r/2}\Delta^{-1/2}|t|^{-r/2}\},
 \qquad B(U)\leq e^{r/2}.
\end{equation}
If \(r>2\), then \(Y\) has a continuous density vanishing at infinity,
and
\begin{equation}\label{four:diffuse-density}
 \|p_Y\|_\infty
 \leq\frac r{r-2}B(U)^{2/r}\Delta^{-1/r}
 \leq\frac{er}{r-2}\Delta^{-1/r}.
\end{equation}
The same conclusion holds if the indicated representation is available
only for a principal block in the original coordinates of a quadratic
form; all other blocks and couplings, as well as affine terms, may be
arbitrary.
\end{theorem}
\begin{proof}
For fixed \(t\), let \(h_{i,t}(x)=f_i(x)e^{itV_i(x)}\).  Push the
complex product measure \(\prod_i h_{i,t}(x_i)\,dx_i\) forward by
\(x\mapsto U^{\mathsf T}x\).  Integrating over the orthogonal
complement, as in Lemma~\ref{four:projection}, gives an \(L^1\)
density \(g_t\) with \(\|g_t\|_1\leq1\).  Fubini gives
\[
 \widehat g_t(\xi)
 =\prod_i\widehat h_{i,t}(\langle U^{\mathsf T}e_i,\xi\rangle).
\]
Rows with \(c_i=0\) contribute factors of modulus at most one.
For the other rows set \(v_i=U^{\mathsf T}e_i/\sqrt{c_i}\).
Then \(\sum_i c_iv_iv_i^{\mathsf T}=I_r\).
The geometric Brascamp--Lieb inequality applied to
\(H_i(s)=|\widehat h_{i,t}(\sqrt{c_i}s)|^{1/c_i}\) yields
\[
 \|\widehat g_t\|_1
 \leq\prod_{i:c_i>0}
          c_i^{-c_i/2}\|\widehat h_{i,t}\|_{1/c_i}.
\]
The output exponent \(1/c_i\) is at least two.  Sharp
Hausdorff--Young with conjugate exponent \(1/(1-c_i)\) gives
\[
 \|\widehat h_{i,t}\|_{1/c_i}
 \leq(2\pi)^{c_i}
       (1-c_i)^{-(1-c_i)/2}c_i^{c_i/2}.
\]
The \(c_i\) factors cancel, and \(\sum_i c_i=r\), so
\begin{equation}\label{four:complex-projection}
 \|\widehat g_t\|_1\leq(2\pi)^rB(U).
\end{equation}
The entropy inequality \eqref{four:entropy} gives \(B(U)\leq e^{r/2}\).

The characteristic function is exactly
\(\int e^{it y^{\mathsf T}\Lambda y}g_t(y)\,dy\).
Lemma~\ref{four:fresnel} and \eqref{four:complex-projection} prove
\eqref{four:diffuse-decay}; the bound one follows from the probability
law of \(Y\).  Lemma~\ref{four:inversion}, with
\(b=r/2\) and \(a=\pi B(U)^{2/r}\Delta^{-1/r}\), proves
\eqref{four:diffuse-density}.

For the principal-block statement, condition on the complementary
original coordinates.  The remaining coordinates keep their original
independent distributions.  The outside couplings induce only linear
terms in them, and these are absorbed in the functions \(V_i\).
The bound \eqref{four:diffuse-decay} is uniform in all conditioning
values, so averaging preserves it and its integrability.
\end{proof}

\subsection{Decoupling and volume-sampling Brascamp--Lieb inequalities}

\begin{lemma}[A cross-block decoupling inequality]\label{four:decoupling}
Let \(I,J\) be disjoint original coordinate sets, let \(M\) be a real
\(|I|\times|J|\) matrix, and let
\[
 Y=F(X_I)+2X_I^{\mathsf T}MX_J+G(X_J),
\]
where \(F,G\) are arbitrary real measurable functions finite almost
surely.  If \(X'_J\) is an independent copy and \(D=X_J-X'_J\), then
\begin{equation}\label{four:decoupling-bound}
 |\varphi_Y(t)|^2
 \leq\mathbb E\prod_{i\in I}
      |\widehat f_i(2t(MD)_i)|.
\end{equation}
For every \(q\geq1\), consequently,
\begin{equation}\label{four:decoupling-q}
 |\varphi_Y(t)|^{2q}
 \leq\mathbb E\prod_{i\in I}
      |\widehat f_i(2t(MD)_i)|^q .
\end{equation}
Each coordinate of \(D\) has its own difference density and these
coordinates are independent.
\end{lemma}
\begin{proof}
Write the characteristic function by first integrating \(X_J\).
Cauchy--Schwarz in the probability law of \(X_I\), after taking
absolute values of the factor \(e^{itF(X_I)}\), gives
\[
 |\varphi_Y(t)|^2
 \leq\mathbb E_{X_I}
 \left|\mathbb E_{X_J}
   e^{itG(X_J)+2itX_I^{\mathsf T}MX_J}\right|^2.
\]
Expand the square using \(X_J'\).  Integration in \(X_I\) gives the
product of the scalar transforms in \eqref{four:decoupling-bound},
multiplied by \(e^{it(G(X_J)-G(X'_J))}\).
The triangle inequality proves \eqref{four:decoupling-bound}.
Jensen's inequality for the probability law of \(D\) proves
\eqref{four:decoupling-q}.  Independence and the difference-density
description follow directly from the independent coordinate pairs.
\end{proof}

For a real \(m\times d\) matrix \(M\) of rank at least \(k\), write
\begin{equation}\label{four:volume-definition}
 Z_k=e_k(M^{\mathsf T}M)
 =\sum_{\substack{S\subseteq[m],\,T\subseteq[d]\\|S|=|T|=k}}
           \det(M_{S,T})^2>0.
\end{equation}
Here \(e_k\) is the elementary symmetric polynomial of the squared
singular values.  The equality follows from Cauchy--Binet.
The squared-minor probability law and its marginals are
\begin{equation}\label{four:volume-marginals}
 w_{S,T}=\frac{\det(M_{S,T})^2}{Z_k},\qquad
 \alpha_i=\sum_{S\ni i,T}w_{S,T},\qquad
 \beta_j=\sum_{S,T\ni j}w_{S,T}.
\end{equation}
Thus \(0\leq\alpha_i,\beta_j\leq1\) and
\(\sum_i\alpha_i=\sum_j\beta_j=k\).

\begin{lemma}[Volume-sampling Fourier estimate]\label{four:volume}
In the setting of Lemma~\ref{four:decoupling}, suppose \(Z_k>0\).
Let \(q\geq1\) and assume
\begin{equation}\label{four:volume-gate}
 \alpha_i\leq q/2\quad\hbox{for every row }i.
\end{equation}
Then, for \(t\ne0\),
\begin{align}
 |\varphi_Y(t)|^{2q}
 &\leq e^k q^{-k/2}
          \left(\frac{\pi}{|t|}\right)^kZ_k^{-1/2},
 \label{four:volume-sharp-decay}\\
 |\varphi_Y(t)|^{2q}
 &\leq e^{k/2}q^{-k/2}
          \left(\frac{2\pi}{|t|}\right)^kZ_k^{-1/2}.
 \label{four:volume-classical-decay}
\end{align}
The first bound is stronger.  Neither bound requires a condition on
the column marginals.
\end{lemma}
\begin{proof}
Put \(h_i(u)=|\widehat f_i(2tu)|^q\), and let \(g_j\) be the
difference density of \(D_j\).  By \eqref{four:decoupling-q} it
suffices to bound
\[
 \int_{\mathbb R^d}
        \prod_i h_i((Mx)_i)\prod_jg_j(x_j)\,dx.
\]
Use row maps \(x\mapsto(Mx)_i\) with exponents \(\alpha_i\),
and coordinate maps \(x\mapsto x_j\) with exponents
\(1-\beta_j\).  Their sum is \(d\).  Rows with \(\alpha_i=0\)
can be dropped because \(h_i\leq1\).  Coordinates with
\(\beta_j=1\) can be dropped because \(g_j\leq1\).
Every retained row map is nonzero: a zero row occurs in no
nonzero minor.

We first bound the Brascamp--Lieb constant of this datum.  Let
\(a_i>0\) be Gaussian row parameters after absorbing the
\(\alpha_i\) factors, and let \(b_j>0\) be the similarly absorbed
parameters of the active coordinate maps.  Set \(a_i=0\) for
omitted rows and \(b_j=0\) when \(\beta_j=1\).
Expansion of a diagonal-plus-positive-semidefinite determinant,
followed by Cauchy--Binet, gives
\begin{align}
 \det(D_b+M^{\mathsf T}D_aM)
 &=\sum_{T\subseteq[d]}
       \left(\prod_{j\notin T}b_j\right)
       \det\big((M^{\mathsf T}D_aM)_{T,T}\big)
 \nonumber\\
 &=\sum_{\substack{S\subseteq[m],\,T\subseteq[d]\\|S|=|T|}}
       \det(M_{S,T})^2
       \prod_{i\in S}a_i\prod_{j\notin T}b_j .
 \label{four:determinant-expansion}
\end{align}
This is a polynomial identity and is valid when some diagonal
parameters vanish.  Retain the terms with \(|S|=|T|=k\).
Weighted arithmetic--geometric mean with probabilities
\(w_{S,T}\) gives
\begin{equation}\label{four:determinant-amgm}
 \det(D_b+M^{\mathsf T}D_aM)
 \geq Z_k\prod_i a_i^{\alpha_i}
                    \prod_j b_j^{1-\beta_j}.
\end{equation}
Every column with \(\beta_j=1\) belongs to every \(T\) of positive
weight, so its zero parameter never occurs in the product of an
active summand.  The same observation applies to omitted rows.
Thus all active products in the arithmetic--geometric mean step
are positive, and no inverse of a zero parameter is used.

Substitute \eqref{four:determinant-amgm} in
\eqref{four:gaussian-bl}, restoring the absorbed exponents.
The Brascamp--Lieb constant is at most
\begin{equation}\label{four:volume-bl-constant}
 Z_k^{-1/2}\prod_i\alpha_i^{-\alpha_i/2}
                 \prod_j(1-\beta_j)^{-(1-\beta_j)/2}.
\end{equation}
Zero exponents in this formula have value one by continuity.
Consequently the integral under consideration is at most this
constant times
\(\prod_i\|h_i\|_{1/\alpha_i}
  \prod_j\|g_j\|_{1/(1-\beta_j)}\).
The latter column norms are at most one.

The scalar estimate \eqref{four:fourier-power-norm}, including the
change of variable \(2tu\), gives
\begin{equation}\label{four:volume-hy}
 \|h_i\|_{1/\alpha_i}
 \leq\left(\frac{\pi}{|t|}\right)^{\alpha_i}
       \left(\frac{\alpha_i}{q}\right)^{\alpha_i/2}
       \left(1-\frac{\alpha_i}{q}\right)^{-(q-\alpha_i)/2}.
\end{equation}
Its use is legitimate precisely because \(q/\alpha_i\geq2\).
After cancellation of the powers of \(\alpha_i\), the row
entropy factor is at most \(e^{k/2}\), by
\eqref{four:entropy} applied to \(x=\alpha_i/q\).
The column entropy factor in \eqref{four:volume-bl-constant}
is also at most \(e^{k/2}\).  Since \(\sum_i\alpha_i=k\),
this proves \eqref{four:volume-sharp-decay}.

Alternatively, use only the last bound in
\eqref{four:lp-sharp} to obtain
\[
 \|h_i\|_{1/\alpha_i}
 \leq\left(\frac{2\pi}{|t|}\right)^{\alpha_i}
       \left(\frac{\alpha_i}{q}\right)^{\alpha_i/2}.
\]
The same cancellation, followed by just the column entropy
bound, proves \eqref{four:volume-classical-decay}.
The stronger estimate follows also numerically from
\(e<2\sqrt e\).
\end{proof}

\begin{theorem}[General volume density criterion]\label{four:volume-density}
Under the hypotheses of Lemma~\ref{four:volume}, suppose in addition
that \(k>2q\).  Then \(Y\) has a continuous density vanishing at
infinity and
\begin{equation}\label{four:volume-density-bound}
 \|p_Y\|_\infty
 \leq\frac{ke}{(k-2q)\sqrt q}\,Z_k^{-1/(2k)}.
\end{equation}
Using the second estimate of that lemma gives the also valid bound
\begin{equation}\label{four:volume-density-classical}
 \|p_Y\|_\infty
 \leq\frac{2k\sqrt e}{(k-2q)\sqrt q}\,Z_k^{-1/(2k)}.
\end{equation}
Both conclusions apply to a cross block \(M=A_{I,J}\) in a larger
quadratic form with arbitrary self blocks, outside couplings, and
affine terms.
\end{theorem}
\begin{proof}
The first Fourier envelope has exponent \(b=k/(2q)>1\).
Its crossover scale is
\(a=e\pi q^{-1/2}Z_k^{-1/(2k)}\).
Lemma~\ref{four:inversion} proves
\eqref{four:volume-density-bound}.  The second crossover scale is
\(2\pi\sqrt e\,q^{-1/2}Z_k^{-1/(2k)}\), giving
\eqref{four:volume-density-classical}.
For a larger form, condition on the coordinates outside \(I\cup J\).
They change only the functions \(F,G\) and an additive constant.
The Fourier envelopes remain uniform and can therefore be averaged.
\end{proof}

\begin{corollary}[Three-volume and five-volume cases]\label{four:d-cases}
The following explicit cases of Theorem~\ref{four:volume-density}
hold:
\begin{enumerate}
\item If \(Z_3>0\) and \(\max_i\alpha_i^{(3)}\leq1/2\), then
\[
 \|p_Y\|_\infty\leq3e\,Z_3^{-1/6}
                         \leq6\sqrt e\,Z_3^{-1/6}.
\]
\item If \(Z_3>0\) and
\(\max_i\alpha_i^{(3)}\leq h\) with \(1/2\leq h<3/4\), then
\[
 \|p_Y\|_\infty
 \leq\frac{3e}{(3-4h)\sqrt{2h}}\,Z_3^{-1/6}
 \leq\frac{6\sqrt e}{(3-4h)\sqrt{2h}}\,Z_3^{-1/6}.
\]
\item If \(Z_5>0\), with no marginal hypothesis, then
\[
 \|p_Y\|_\infty\leq\frac{5e}{\sqrt2}\,Z_5^{-1/10}
                         \leq5\sqrt{2e}\,Z_5^{-1/10}.
\]
\end{enumerate}
These are respectively the three-volume, fractional three-volume,
and unconditional five-volume criteria.
\end{corollary}
\begin{proof}
Take \((k,q)=(3,1),(3,2h),(5,2)\), respectively.  In the last
case every inclusion probability is at most one, so the gate
\(\alpha_i^{(5)}\leq q/2\) is automatic.  The strict inequality
\(k>2q\) holds in all three cases.
\end{proof}

\begin{remark}[Choice of a minor law]\label{four:weighted-volume}
The canonical squared-minor probabilities can be replaced by any
probability law \(v_{S,T}\) supported on nonzero \(k\times k\)
minors.  If its row marginals satisfy the same gate, the bound
\eqref{four:volume-density-bound} acquires the factor
\[
 \exp\!\left(\frac{D(v\Vert w)}{2k}\right),\qquad
 D(v\Vert w)=\sum_{S,T}v_{S,T}\log\frac{v_{S,T}}{w_{S,T}}.
\]
Indeed, apply weighted arithmetic--geometric mean to
\(\sum v_{S,T}
 [\det(M_{S,T})^2/v_{S,T}]
 \prod_{i\in S}a_i\prod_{j\notin T}b_j\).
The determinant lower bound has \(Z_k e^{-D(v\Vert w)}\) in
place of \(Z_k\).  All subsequent steps are identical.  This
form allows the marginal constraint to be traded against an
explicit entropy cost.
\end{remark}

\subsection{Exact marginal factors and entropy optimization}

\begin{theorem}[Exact weighted-volume factors]\label{four:exact-marginal}
Let \(w_{S,T}\) be any probability law on the nonzero
\(k\times k\) minors of a physical cross block \(M\), and let
\(\alpha_i,\beta_j\) be its row and column inclusion probabilities.
Set
\[
 G(w)=\exp\left(\sum_{S,T}w_{S,T}
                 \log\frac{\det(M_{S,T})^2}{w_{S,T}}\right).
\]
For \(q\geq1\) with \(\alpha_i\leq q/2\), define
\[
 R_q(w)=
 \prod_i\left(1-\frac{\alpha_i}{q}\right)^{-(q-\alpha_i)/2}
 \prod_j(1-\beta_j)^{-(1-\beta_j)/2}.
\]
All zero-exponent factors have their continuous value one.
Then, in the setting of Lemma~\ref{four:decoupling},
\begin{equation}\label{four:exact-marginal-decay}
 |\varphi_Y(t)|^{2q}
 \leq R_q(w)q^{-k/2}\pi^kG(w)^{-1/2}|t|^{-k}.
\end{equation}
If \(k>2q\), it follows that
\begin{equation}\label{four:exact-marginal-density}
 \|p_Y\|_\infty
 \leq\frac{k}{(k-2q)\sqrt q}\,
                  R_q(w)^{1/k}G(w)^{-1/(2k)}.
\end{equation}
The density is continuous and vanishes at infinity.
In particular, for \(k=3,q=1\), writing \(R=R_1\),
\begin{equation}\label{four:q1-exact}
 \|p_Y\|_\infty\leq3R(w)^{1/3}G(w)^{-1/6},\qquad
 \|\varphi_Y\|_1\leq6\pi R(w)^{1/3}G(w)^{-1/6}.
\end{equation}
\end{theorem}
\begin{proof}
In the determinant expansion
\eqref{four:determinant-expansion}, apply arithmetic--geometric
mean with the law \(w\) to obtain
\[
 \det(D_b+M^{\mathsf T}D_aM)
 \geq G(w)\prod_i a_i^{\alpha_i}
                   \prod_jb_j^{1-\beta_j}.
\]
The argument for saturated marginals in the proof of
Lemma~\ref{four:volume} applies without change.
Thus its Brascamp--Lieb constant is
\eqref{four:volume-bl-constant} with \(G(w)\) in place of \(Z_k\).
Use the exact scalar estimate \eqref{four:volume-hy} and leave
the row and column factors unrelaxed.
The powers of \(\alpha_i\) cancel, and
\(\sum_i\alpha_i=k\); the result is
\eqref{four:exact-marginal-decay}.
Its exponent is \(k/(2q)\) and its crossover scale is
\(\pi q^{-1/2}R_q(w)^{1/k}G(w)^{-1/(2k)}\).
Lemma~\ref{four:inversion} proves
\eqref{four:exact-marginal-density} and the last assertion.
\end{proof}

The entropy inequality gives \(R_q(w)\leq e^k\), recovering
the coefficient in Lemma~\ref{four:volume}.  For the canonical squared-minor law
\(p_{S,T}=\det(M_{S,T})^2/Z_k\), the identity
\[
 G(w)=Z_k e^{-D(w\Vert p)}
\]
shows that \(p\) maximizes \(G\) without marginal constraints.
It does not assert that \(p\) minimizes the whole expression
in \eqref{four:exact-marginal-density}; the marginal factor
depends on the law as well.

\begin{proposition}[Balanced six-row and six-column blocks]
\label{four:balanced-six}
If a law on \(3\times3\) minors of a six-row, six-column
physical cross block has every row and column marginal \(1/2\),
then
\[
 \|p_Y\|_\infty\leq6G(w)^{-1/6}.
\]
For example, let \(u=6^{-1/2}(1,\ldots,1)^{\mathsf T}\) and
\[
 M=sI_6+(L-s)uu^{\mathsf T},\qquad L>s>0.
\]
With arbitrary self functions on the two banks,
\begin{equation}\label{four:balanced-six-bound}
 \|p_Y\|_\infty
 \leq6[10s^4(L^2+s^2)]^{-1/6}.
\end{equation}
This example need not satisfy the sufficient projection condition
of Corollary~\ref{four:tail}.
\end{proposition}
\begin{proof}
Each of the twelve marginal factors equals \(2^{1/4}\),
so \(R(w)=8\).  Formula \eqref{four:q1-exact} proves the
first assertion.
The displayed matrix has singular values \(L,s,s,s,s,s\).
Simultaneous coordinate-permutation symmetry makes its canonical
row and column marginal vectors constant; since each sums to
three, all entries are \(1/2\).  Moreover
\[
 Z_3=\binom52 L^2s^4+\binom53s^6
     =10s^4(L^2+s^2).
\]
The canonical law has \(G=Z_3\), proving the bound.
For the final assertion, any projection with all diagonal entries
at most \(1/2\) has rank at most three and is therefore proper.
For a unit \(v\) in its orthogonal complement,
\(\|M^{\mathsf T}v\|\geq s\); hence
\(\|(I-P)M\|_{\rm op}\geq s\).
Here \(T_3=4s^2\), so the sufficient requirement
\(\|(I-P)M\|\leq\sqrt{T_3}/15=2s/15\) fails.
\end{proof}

We next optimize the determinant entropy subject to a row
marginal constraint.  All optimization sets are finite
dimensional, and all their coordinates refer to original rows.
For a row \(k\)-set \(S\), define
\[
 a_S=\det(M_SM_S^{\mathsf T})
     =\sum_{|T|=k}\det(M_{S,T})^2,\qquad
 \mathcal B=\{S:a_S>0\},\qquad Z_k=\sum_{S\in\mathcal B}a_S.
\]
For a probability law \(\mu\) on \(\mathcal B\), let
\[
 \alpha_i(\mu)=\sum_{S\ni i}\mu_S,\qquad
 G(\mu)=\exp\left(\sum_S\mu_S\log(a_S/\mu_S)\right).
\]
Terms with \(\mu_S=0\) are interpreted by continuity.

\begin{lemma}[Reduction to row entropy]\label{four:row-entropy}
For fixed row law \(\mu\), the lift
\begin{equation}\label{four:canonical-lift}
 w_{S,T}=\mu_S\frac{\det(M_{S,T})^2}{a_S}
\end{equation}
has the same row marginals and maximizes \(G(w)\) among all
minor laws with row law \(\mu\).  Its value is \(G(\mu)\).
If \(p_S=a_S/Z_k\), then
\begin{equation}\label{four:row-kl}
 G(\mu)=Z_k e^{-D(\mu\Vert p)}.
\end{equation}
\end{lemma}
\begin{proof}
For a general lift let \(\nu_S\) be the conditional column
law given \(S\), and let
\(\pi_S(T)=\det(M_{S,T})^2/a_S\).
Expanding logarithms gives
\[
 \log G(w)
 =\sum_S\mu_S\log(a_S/\mu_S)
                  -\sum_S\mu_SD(\nu_S\Vert\pi_S).
\]
Relative entropy is nonnegative, as follows from
\(\log x\leq x-1\), with equality for identical laws.
Thus the maximum is attained by \(\nu_S=\pi_S\), which is
\eqref{four:canonical-lift}.  Expanding
\(D(\mu\Vert p)\) proves \eqref{four:row-kl}.
\end{proof}

\begin{theorem}[Gibbs dual and concentration on a bounded row set]
\label{four:gibbs}
Assume \(Z_k>0\).  For \(0<h<1\), let
\[
 \mathcal F_h=\{\mu:\alpha_i(\mu)\leq h\text{ for every }i\},
 \qquad G_h(M)=\max_{\mu\in\mathcal F_h}G(\mu),
\]
with \(G_h(M)=0\) if \(\mathcal F_h\) is empty.  Then
\begin{equation}\label{four:gibbs-dual}
 G_h(M)=\inf_{\lambda_i\geq0}
  \exp\left(h\sum_i\lambda_i\right)
       \sum_{S\in\mathcal B}a_S
                    \exp\left(-\sum_{i\in S}\lambda_i\right).
\end{equation}
If a finite vector \(\lambda\) attains the infimum, the
corresponding maximizing law is
\begin{equation}\label{four:gibbs-law}
 \mu_\lambda(S)=
 \frac{a_S e^{-\sum_{i\in S}\lambda_i}}
      {\sum_Ra_R e^{-\sum_{i\in R}\lambda_i}},
\end{equation}
and \(\lambda_i>0\) implies \(\alpha_i(\mu_\lambda)=h\).
Regardless of dual attainment, there exists a set \(H\) of rows
such that
\begin{equation}\label{four:gibbs-delete}
 |H|<k/h,\qquad
 e_k(M_{[m]\setminus H}^{\mathsf T}M_{[m]\setminus H})
                         \leq G_h(M).
\end{equation}
\end{theorem}
\begin{proof}
Let \(\Delta\) be the probability simplex on \(\mathcal B\),
and define its continuous concave entropy function
\[
 F(\mu)=\sum_S\mu_S\log(a_S/\mu_S).
\]
For \(\lambda_i\geq0\), put
\[
 \mathcal L(\mu,\lambda)
 =F(\mu)+\sum_i\lambda_i(h-\alpha_i(\mu)).
\]
The Gibbs variational identity states
\begin{equation}\label{four:gibbs-variational}
 \max_{\mu\in\Delta}\mathcal L(\mu,\lambda)
 =h\sum_i\lambda_i+
         \log\sum_S a_S e^{-\sum_{i\in S}\lambda_i}.
\end{equation}
Indeed, subtracting the right side from
\(\mathcal L(\mu,\lambda)\) gives
\(-D(\mu\Vert\mu_\lambda)\).

We justify the dual equality also when constraints are attained
only on the boundary or are infeasible.
Restrict first to \(0\leq\lambda_i\leq N\).
The compact convex minimax theorem \citep{Sion1958}, applied to
the continuous concave-affine function \(\mathcal L\), gives
\[
 \min_{\lambda\in[0,N]^m}\max_{\mu\in\Delta}\mathcal L(\mu,\lambda)
 =\max_{\mu\in\Delta}
   \left[F(\mu)-N\sum_i(\alpha_i(\mu)-h)_+\right].
\]
If \(\mathcal F_h\ne\varnothing\), the right side decreases,
as \(N\) increases, to \(\max_{\mathcal F_h}F\).
To see this explicitly, take maximizing laws along a sequence
\(N\to\infty\) and a convergent subsequence in the compact
simplex.  Since \(F\) is bounded and a feasible law gives a
fixed finite lower bound, their constraint violations tend
to zero.  The limiting law is feasible.  Dropping the
nonpositive penalty and using continuity of \(F\) proves that
the limiting value is at most the constrained maximum; the
reverse inequality holds for every \(N\).
If \(\mathcal F_h\) is empty, the continuous sum of positive
constraint violations has a strictly positive minimum on
\(\Delta\), so these maxima tend to \(-\infty\).
Taking \(N\to\infty\) and using
\eqref{four:gibbs-variational} proves the logarithm of
\eqref{four:gibbs-dual}, including the value \(-\infty\)
in the infeasible case.

If the dual minimizer is finite, differentiation of its
logarithmic objective gives \(h-\alpha_i(\mu_\lambda)\).
Optimality on the nonnegative orthant therefore implies
\(\alpha_i(\mu_\lambda)\leq h\), and equality if
\(\lambda_i>0\).  Substitution in \(\mathcal L\) shows that
this Gibbs law attains the primal maximum.

It remains to prove \eqref{four:gibbs-delete}.
If \(G_h=Z_k\), take \(H=\varnothing\).
Otherwise choose, for arbitrarily small \(\varepsilon>0\),
a nonnegative vector \(\lambda\) whose dual objective is
at most \(G_h+\varepsilon<Z_k\).
Then \(\Lambda=\sum_i\lambda_i>0\).
Put \(\omega_i=\lambda_i/\Lambda\) and
\[
 H=\{i:\omega_i>h/k\}.
\]
Since \(\sum_i\omega_i=1\), necessarily \(|H|<k/h\).
For every \(k\)-set \(S\) avoiding \(H\),
\(\sum_{i\in S}\lambda_i\leq h\Lambda\).
The corresponding term in the dual objective is therefore
at least \(a_S\).  Thus
\[
 \sum_{S:S\cap H=\varnothing}a_S\leq G_h+\varepsilon.
\]
The left side is the residual \(k\)-volume in
\eqref{four:gibbs-delete}.  There are only finitely many
possible sets \(H\); along a sequence
\(\varepsilon\downarrow0\), one repeats infinitely often.
Passing to that subsequence proves the desired bound,
including when \(G_h=0\).
\end{proof}

\begin{corollary}[A three-volume or bounded-row alternative]
\label{four:gibbs-alternative}
Let \(1/2\leq h<3/4\) and \(g>0\).
For a physical cross block \(M\) with \(Z_3>0\), either
\[
 \|p_Y\|_\infty
 \leq\frac{3e}{(3-4h)\sqrt{2h}}\,g^{-1/6},
\]
with a continuous density vanishing at infinity, or there is
a row set \(H\) with \(|H|<3/h\) for which the residual
third volume is less than \(g\).
In particular \(h=3/5\) requires at most four rows in the
second alternative.
\end{corollary}
\begin{proof}
If \(G_h\geq g\), compactness gives an optimizing feasible law.
Lift it by Lemma~\ref{four:row-entropy} and apply
Theorem~\ref{four:exact-marginal} with \(q=2h\) and
\(R_q\leq e^3\).  If \(G_h<g\), use
\eqref{four:gibbs-delete}.  At \(h=3/5\), the strict bound
\(|H|<5\) is equivalent to \(|H|\leq4\).
\end{proof}

\begin{proposition}[A single active row constraint]
\label{four:one-hub-entropy}
Fix a distinguished row \(i_0\), and write
\[
 Z_1=\sum_{S\ni i_0}a_S,\qquad
 Z_0=\sum_{S\not\ni i_0}a_S.
\]
Suppose \(Z_0,Z_1>0\), \(p=Z_1/(Z_0+Z_1)>h\), and the law
\[
 \mu(S)=
 \begin{cases}
  h\,a_S/Z_1,&i_0\in S,\\
  (1-h)a_S/Z_0,&i_0\notin S
 \end{cases}
\]
satisfies all other row constraints.  Then
\begin{equation}\label{four:one-hub-formula}
 G_h(M)=\left(\frac{Z_1}{h}\right)^h
         \left(\frac{Z_0}{1-h}\right)^{1-h}
       =(Z_0+Z_1)e^{-D_{\rm Ber}(h\Vert p)}.
\end{equation}
A dual optimizer has all multipliers zero except
\[
 e^{-\lambda_{i_0}}=\frac{hZ_0}{(1-h)Z_1}.
\]
\end{proposition}
\begin{proof}
If a row law gives total mass \(x\) to sets containing \(i_0\),
the entropy chain rule bounds its logarithmic value by
\[
 x\log(Z_1/x)+(1-x)\log(Z_0/(1-x)),
\]
with equality exactly when the conditional laws in the two
classes are canonical.  The derivative of this expression is
\(\log[Z_1(1-x)/(Z_0x)]\), which is positive for \(x<p\).
Under \(x\leq h<p\), its maximum is therefore attained at \(h\).
The displayed law attains this value and, by assumption, is
feasible for every other constraint.  This proves
\eqref{four:one-hub-formula}.  Substitution in
\eqref{four:gibbs-law} verifies the stated multiplier and
complementary slackness.
\end{proof}

\subsection{Regularized determinants and ridge leverage}

\begin{lemma}[A graph inequality with a free regularizer]
\label{four:graph}
Let \(R\) be a real \(m\times k\) matrix with columns \(r_j\).
Let \(D_1,\ldots,D_m\) be independent differences of independent
copies of variables with densities bounded by one, and let
\(f_1,\ldots,f_k\) be further probability densities bounded by one.
Fix \(\alpha>0\), \(a\in\mathbb R\), and \(1\leq q\leq2\).  Put
\[
 H=\alpha I_m+a^2RR^{\mathsf T},\qquad
 c_i=\alpha e_i^{\mathsf T}H^{-1}e_i,\qquad
 d_j=a^2r_j^{\mathsf T}H^{-1}r_j,\qquad
 \theta=\sum_jd_j=\sum_i(1-c_i).
\]
If \(d_j\leq q/2\) for every \(j\), then, with
\(L_q=\pi e\sqrt{2/q}\),
\begin{equation}\label{four:graph-bound}
 \mathbb E\prod_j|\widehat f_j(a\langle r_j,D\rangle)|^q
 \leq
 \det(I_m+a^2RR^{\mathsf T}/\alpha)^{-1/2}
          (L_q/\sqrt\alpha)^\theta .
\end{equation}
\end{lemma}
\begin{proof}
For \(a=0\), both the expectation and the determinant are one,
and \(\theta=0\).  Assume \(a\ne0\).
The vectors
\[
 u_i=\sqrt\alpha H^{-1/2}e_i,\qquad
 v_j=aH^{-1/2}r_j
\]
satisfy
\(\sum_i u_iu_i^{\mathsf T}+\sum_jv_jv_j^{\mathsf T}=I_m\).
Their squared lengths are \(c_i,d_j\), respectively; thus
\(0<c_i\leq1\), \(0\leq d_j\leq1\).
The trace identity gives the displayed expression for \(\theta\).
Zero \(v_j\) may be omitted.

Let \(g_i\) be the density of \(D_i\), and put
\(h_j=|\widehat f_j|^q\).  Change variables \(x=H^{-1/2}z\)
in the expectation integral and normalize the nonzero vectors
\(u_i,v_j\) to unit vectors.  The geometric Brascamp--Lieb
inequality gives
\begin{equation}\label{four:graph-bl}
 \det(H)^{-1/2}
 \prod_i\left[\alpha^{c_i/2}c_i^{-c_i/2}\|g_i\|_{1/c_i}\right]
 \prod_j\left[d_j^{-d_j/2}\|h_j\|_{1/d_j}\right].
\end{equation}
For example, the argument of \(g_i\) is
\(\sqrt{c_i/\alpha}\) times the relevant unit linear form, which
accounts for its factor \(\alpha^{c_i/2}c_i^{-c_i/2}\).

Define
\[
 G(c)=2^c(1+c)^{-(1+c)/2},\qquad
 H_q(d)=q^{q/2-d}(q-d)^{-(q-d)/2}.
\]
Lemma~\ref{four:scalar-norms} cancels the powers of \(c_i,d_j\)
in \eqref{four:graph-bl}, leaving
\begin{equation}\label{four:graph-after-scalar}
 \det(H)^{-1/2}\alpha^{(m-\theta)/2}(2\pi)^\theta
                     \prod_iG(c_i)\prod_jH_q(d_j).
\end{equation}
The logarithms of these scalar functions are concave:
\[
 (\log G)''(c)=-\frac1{2(1+c)},\qquad
 (\log H_q)''(d)=-\frac1{2(q-d)}.
\]
The tangent to \(\log G\) at \(c=1\), and the tangent to
\(\log H_q\) at \(d=0\), therefore give
\[
 \log G(c)\leq\frac{1-\log2}{2}(1-c),\qquad
 \log H_q(d)\leq\frac{1-\log q}{2}d.
\]
Their endpoint values are \(G(1)=H_q(0)=1\).
The product in \eqref{four:graph-after-scalar} is consequently
at most \(\exp\{[1-\tfrac12\log(2q)]\theta\}\).
Factor \(\alpha^m\) from \(\det H\).  The remaining scalar
multiplier is
\[
 \left(\frac{2\pi e}{\sqrt{2q\alpha}}\right)^\theta
 =(L_q/\sqrt\alpha)^\theta,
\]
which proves \eqref{four:graph-bound}.
\end{proof}

\begin{theorem}[Regularized cross-block determinant]
\label{four:regularized}
Let \(Y\) have the form in Lemma~\ref{four:decoupling}, with cross
block \(B\).  Let \(1\leq q\leq2\).  Suppose that either \(q=2\)
or the projection \(P_{\operatorname{range}(B)}\), acting on the
row-index space, satisfies
\[
 \max_i(P_{\operatorname{range}(B)})_{ii}\leq q/2.
\]
Then
\begin{equation}\label{four:regularized-bound}
 |\varphi_Y(t)|
 \leq\det\left(I+\frac{2qt^2}{\pi^2e^2}BB^{\mathsf T}\right)^{-1/(4q)}.
\end{equation}
In particular, without any leverage hypothesis,
\begin{equation}\label{four:unrestricted-determinant}
 |\varphi_Y(t)|
 \leq\det\left(I+\frac{4t^2}{\pi^2e^2}BB^{\mathsf T}\right)^{-1/8}.
\end{equation}
If \(B\) has at least \(r>2q\) singular values at least \(\sigma>0\),
the indicated hypothesis implies a continuous density with
\begin{align}
 \|p_Y\|_\infty
 &\leq
 \frac{e\sqrt\pi}{2\sqrt{2q}\,\sigma}
 \frac{\Gamma(r/(4q)-1/2)}{\Gamma(r/(4q))}
 \label{four:regularized-gamma}\\
 &\leq\frac{e}{\sqrt{2q}\,\sigma}\frac r{r-2q}.
 \label{four:regularized-simple}
\end{align}
All these statements allow arbitrary functions in the two self blocks,
and conditioning on additional original coordinates.
\end{theorem}
\begin{proof}
Apply \eqref{four:decoupling-q} and then Lemma~\ref{four:graph}
with \(R=B^{\mathsf T}\), \(a=2t\), and
\(\alpha=L_q^2\).  The graph weights on the Fourier factors are
the diagonal entries of
\[
 (2t)^2B(L_q^2I+(2t)^2B^{\mathsf T}B)^{-1}B^{\mathsf T}.
\]
This matrix is bounded between zero and
\(P_{\operatorname{range}(B)}\), as is seen by singular-value
decomposition.  Its diagonal entries are always at most one,
so \(q=2\) needs no assumption.  The multiplier in
\eqref{four:graph-bound} is one for this choice of \(\alpha\).
Taking the \(2q\)-th root and using equality of the nonzero spectra
of \(BB^{\mathsf T}\) and \(B^{\mathsf T}B\) proves
\eqref{four:regularized-bound}.

Discarding determinant factors other than the \(r\) indicated ones
gives
\[
 |\varphi_Y(t)|
 \leq\left(1+\frac{2q\sigma^2t^2}{\pi^2e^2}\right)^{-r/(4q)}.
\]
The integral
\[
 \int_{\mathbb R}(1+x^2)^{-a}\,dx
 =\sqrt\pi\,\frac{\Gamma(a-1/2)}{\Gamma(a)},\qquad a>1/2,
\]
follows from \(u=x^2\) and the beta integral.  Fourier inversion
now proves \eqref{four:regularized-gamma}.
Bounding the same envelope by
\(\min\{1,(\pi e/(\sqrt{2q}\sigma|t|))^{r/(2q)}\}\)
and applying Lemma~\ref{four:inversion} proves the simpler bound.
The uniformity in self functions and outside coordinates is
inherited from Lemma~\ref{four:decoupling}.
\end{proof}

\begin{theorem}[Ridge-leverage criterion]\label{four:ridge}
For a real cross block \(B\) and \(\gamma>0\), define
\[
 K_\gamma=BB^{\mathsf T}(\gamma^2I+BB^{\mathsf T})^{-1},\quad
 \theta_\gamma=\operatorname{tr}K_\gamma,\quad
 D_\gamma=\det(I+BB^{\mathsf T}/\gamma^2).
\]
Suppose, for some \(1\leq q\leq2\), that
\begin{equation}\label{four:ridge-gate}
 \max_i(K_\gamma)_{ii}\leq q/2,\qquad
 \theta_\gamma>2q.
\end{equation}
Then \(Y\) in Lemma~\ref{four:decoupling} has a continuous density
vanishing at infinity and
\begin{align}
 \|p_Y\|_\infty
 &\leq\frac{e}{\sqrt{2q}\,\gamma}
       \frac{\theta_\gamma}{\theta_\gamma-2q}
       D_\gamma^{-1/(2\theta_\gamma)}
 \label{four:ridge-density}\\
 &\leq\frac{\sqrt e}{\sqrt{2q}\,\gamma}
       \frac{\theta_\gamma}{\theta_\gamma-2q}.
 \label{four:ridge-density-simple}
\end{align}
The two coordinate sets may be interchanged to use the ridge
diagonal on the other side.
\end{theorem}
\begin{proof}
For \(t\ne0\), use the same graph argument as in
Theorem~\ref{four:regularized}, but choose
\(\alpha=4\gamma^2t^2\).  The graph weights \(d_i\) are now
exactly \((K_\gamma)_{ii}\), independent of \(t\); their sum
is \(\theta_\gamma\).  The first condition in
\eqref{four:ridge-gate} makes the scalar Fourier exponents
admissible.  Keeping the multiplier in
\eqref{four:graph-bound}, and taking the \(2q\)-th root, gives
\[
 |\varphi_Y(t)|\leq D_\gamma^{-1/(4q)}
           \left(\frac{L_q}{2\gamma|t|}\right)^{\theta_\gamma/(2q)}.
\]
The power is greater than one.  Its crossover scale is
\[
 a=\frac{L_q}{2\gamma}D_\gamma^{-1/(2\theta_\gamma)}.
\]
Lemma~\ref{four:inversion} and
\(L_q/(2\pi)=e/\sqrt{2q}\) prove
\eqref{four:ridge-density}.
If \(x_j=s_j(B)^2/\gamma^2\), then
\[
 \log D_\gamma=\sum_j\log(1+x_j)
 \geq\sum_j\frac{x_j}{1+x_j}=\theta_\gamma,
\]
because the difference between the summands has derivative
\(x/(1+x)^2\) and vanishes at zero.  Therefore
\(D_\gamma^{-1/(2\theta_\gamma)}\leq e^{-1/2}\), which gives
\eqref{four:ridge-density-simple}.
\end{proof}

\begin{corollary}[Three cross singular values]\label{four:ridge-three}
Suppose \(s_3(B)\geq\sigma>0\), set
\(\gamma=\sigma/\sqrt6\), and assume
\(\max_i(K_\gamma)_{ii}\leq1/2\).  Then
\begin{equation}\label{four:ridge-three-bound}
 \|p_Y\|_\infty\leq\frac{9\sqrt{3e}}{2\sigma}.
\end{equation}
\end{corollary}
\begin{proof}
The first three singular values give
\(\theta_\gamma\geq3/(1+1/6)=18/7\).
Apply \eqref{four:ridge-density-simple} with \(q=1\).
The function \(\theta/(\theta-2)\) decreases for \(\theta>2\),
and its value at \(18/7\) is \(9/2\).
\end{proof}

\begin{lemma}[Stability of ridge leverage]\label{four:ridge-stability}
Let \(B=B_0+E\), \(\|E\|_{\rm op}\leq\eta\), and let \(P\)
be the orthogonal projection onto \(\operatorname{range}(B_0)\).
For every \(\gamma>0\), writing \(\ell_i=P_{ii}\),
\begin{equation}\label{four:ridge-stability-bound}
 (K_\gamma)_{ii}
 \leq\left(\sqrt{\ell_i}
          +\frac{\eta}{\gamma}\sqrt{1-\ell_i}\right)^2.
\end{equation}
In particular, if \(\ell_i\leq\ell\) for all \(i\), the weaker
bound \((\sqrt\ell+\eta/\gamma)^2\) holds.
\end{lemma}
\begin{proof}
The identity
\[
 K_\gamma
 =B(\gamma^2I+B^{\mathsf T}B)^{-1}B^{\mathsf T}
\]
shows \(0\leq K_\gamma\leq I\).
For \(v\perp\operatorname{range}(B_0)\), \(B_0^{\mathsf T}v=0\);
hence
\[
 v^{\mathsf T}K_\gamma v
 \leq\gamma^{-2}\|B^{\mathsf T}v\|^2
 \leq(\eta/\gamma)^2\|v\|^2.
\]
Decompose \(e_i=Pe_i+(I-P)e_i\) and use the triangle inequality
after applying \(K_\gamma^{1/2}\).  Its norm on the first
component is at most \(\sqrt{\ell_i}\), and the preceding
bound on the second is at most
\((\eta/\gamma)\sqrt{1-\ell_i}\).
\end{proof}

\subsection{A spectral-window consequence}

We use the two-part frame partition consequence of
\citet[Corollary 1.5]{MSS2015}: if
\(\sum_i u_iu_i^{\mathsf T}=I_r\) and
\(\|u_i\|^2\leq\varepsilon\), the indices admit a partition
\(I,J\) such that both partial frame operators have norm at most
\((1/\sqrt2+\sqrt\varepsilon)^2\).
The following corollary applies this partition to the range of a
quadratic matrix.

\begin{corollary}[Diffuse full range]\label{four:full-range}
Let \(A\) be symmetric, \(s_3(A)\geq s>0\), and let \(P_A\)
be the projection onto its full range.  If
\(\max_i(P_A)_{ii}\leq1/32\), then, for arbitrary affine terms
and arbitrary additional separable measurable phases,
\begin{equation}\label{four:full-range-bound}
 \|p_{X^{\mathsf T}AX+\sum_iV_i(X_i)}\|_\infty
 \leq\frac{96e}{7\sqrt2\,s}.
\end{equation}
The density is continuous and vanishes at infinity.
\end{corollary}
\begin{proof}
Write \(A=U\Lambda U^{\mathsf T}\), where \(U\) has orthonormal
columns spanning its full range.  The rows of \(U\) form a
tight frame with squared lengths at most \(1/32\).
The cited partition theorem gives
\[
 \frac7{32}I\leq M_I:=U_I^{\mathsf T}U_I,\,
                         M_J:=U_J^{\mathsf T}U_J
                 \leq\frac{25}{32}I.
\]
Indeed \((1/\sqrt2+\sqrt{1/32})^2=25/32\), and the
two operators sum to \(I\).
The nonzero singular values of
\(B=A_{I,J}=U_I\Lambda U_J^{\mathsf T}\) are those of
\(M_I^{1/2}\Lambda M_J^{1/2}\), since
\(U_IM_I^{-1/2}\) and \(U_JM_J^{-1/2}\) have orthonormal
columns.  Multiplication by the two square roots gives
\(s_3(B)\geq(7/32)s\).
Both restricted frames have full column rank, and \(\Lambda\)
is invertible, so \(\operatorname{range}(B)=\operatorname{range}(U_I)\).
Its projection has diagonal
\[
 u_i^{\mathsf T}M_I^{-1}u_i
 \leq\frac{1/32}{7/32}=\frac17<\frac12.
\]
Apply \eqref{four:regularized-simple} with \(q=1\), \(r=3\),
and \(\sigma=7s/32\).
\end{proof}

\begin{theorem}[A diffuse spectral window]\label{four:window}
Let \(A\) be symmetric with \(s_3(A)\geq s>0\), and let \(P\)
be its spectral projection onto the eigenspaces with
\(|\lambda|\geq s/64\).  If
\begin{equation}\label{four:window-hypothesis}
 \max_iP_{ii}\leq1/32,
\end{equation}
then \(X^{\mathsf T}AX+\sum_iV_i(X_i)\), for arbitrary real
measurable \(V_i\) finite almost surely, has a continuous density
vanishing at infinity and
\begin{equation}\label{four:window-bound}
 \|p\|_\infty\leq\frac{288\sqrt{3e}}{13s}<\frac{64}{s}.
\end{equation}
The number, signs, eigenvectors, and nonzero sizes of the smaller
spectral directions are unrestricted.
\end{theorem}
\begin{proof}
Write \(A=A_H+E\), where
\(A_H=U\Lambda_HU^{\mathsf T}\) is the indicated spectral
truncation.  Then \(\|E\|_{\rm op}\leq\eta=s/64\).
Apply the same frame partition as in
Corollary~\ref{four:full-range}.  For
\(B_0=(A_H)_{I,J}\), its proof gives
\[
 s_3(B_0)\geq\frac{7s}{32},\qquad
 \max_i(P_{\operatorname{range}(B_0)})_{ii}\leq\frac17.
\]
The actual cross block is \(B=B_0+E_{I,J}\), and
\(\|E_{I,J}\|_{\rm op}\leq\eta\).  The singular-value
perturbation inequality gives
\[
 s_3(B)\geq\frac{7s}{32}-\frac{s}{64}
          =\sigma:=\frac{13s}{64}.
\]
For completeness, that inequality follows by restricting to the
top-three right-singular subspace of \(B_0\), where
\(\|Bx\|\geq(\frac{7s}{32}-\eta)\|x\|\), and applying the
singular-value min--max principle.

Set \(\gamma=\sigma/\sqrt6\).
Lemma~\ref{four:ridge-stability} yields
\[
 \max_i(K_\gamma)_{ii}
 \leq\left(\frac1{\sqrt7}+\frac{\sqrt6}{13}\right)^2<\frac12.
\]
For example \(1/\sqrt7<2/5\) and \(\sqrt6/13<1/5\), so
the square is less than \(9/25\).
Corollary~\ref{four:ridge-three} proves
\[
 \|p\|_\infty\leq\frac{9\sqrt{3e}}{2\sigma}
              =\frac{288\sqrt{3e}}{13s}.
\]
All small spectral directions have remained in the actual cross
block \(B\); none was removed from the random quadratic expression.
All self terms are permitted by decoupling.  This proves the
claim including arbitrary separable perturbations.
\end{proof}

\subsection{Stability of volume marginals and distributed tails}

\begin{proposition}[Spectral formula for volume marginals]
\label{four:volume-operator}
Let \(M\) have rank at least \(k\), with singular-value
decomposition \(M=U\operatorname{diag}(s_j)V^{\mathsf T}\).
Put \(\lambda_j=s_j^2\), \(Z_k=e_k(\lambda)\), and
\[
 p_j=\frac{\lambda_je_{k-1}(\lambda_{\ne j})}{Z_k},
 \qquad W_k=U\operatorname{diag}(p_j)U^{\mathsf T}.
\]
Then the row marginals in \eqref{four:volume-marginals} satisfy
\((W_k)_{ii}=\alpha_i\), and
\begin{equation}\label{four:volume-operator-bound}
 0\leq W_k\leq I,\qquad
 W_k\leq\frac{e_{k-1}(\lambda)}{Z_k}MM^{\mathsf T}.
\end{equation}
Writing \(T_k=\sum_{j\geq k}\lambda_j>0\), we also have
\begin{equation}\label{four:volume-tail-operator}
 W_k\leq\frac{k}{T_k}MM^{\mathsf T}
       \leq\frac{k}{s_k(M)^2}MM^{\mathsf T}.
\end{equation}
\end{proposition}
\begin{proof}
Multiply row \(i\) of \(M\) by \(\sqrt{1+t}\).
Each squared \(k\)-minor containing that row is multiplied by
\(1+t\), while the other minors are unchanged.  Thus the
coefficient of \(t\) in the new \(Z_k\) is \(Z_k\alpha_i\).
Let \(C=MM^{\mathsf T}\).  Using Sylvester's determinant
identity and the rank-one determinant formula, the generating
polynomial of this row scaling is
\[
 \det(I+zC)
 \left[1+tz\,e_i^{\mathsf T}(I+zC)^{-1}Ce_i\right].
\]
Its coefficient of \(tz^k\), after diagonalizing \(C\), is
\(\sum_jU_{ij}^2\lambda_je_{k-1}(\lambda_{\ne j})\).
Comparison proves \((W_k)_{ii}=\alpha_i\).  This polynomial
argument also covers repeated singular values and zero modes.

The decomposition
\[
 Z_k=\lambda_je_{k-1}(\lambda_{\ne j})
                         +e_k(\lambda_{\ne j})
\]
gives \(0\leq p_j\leq1\), while
\(e_{k-1}(\lambda_{\ne j})\leq e_{k-1}(\lambda)\).
These observations prove \eqref{four:volume-operator-bound}.
Finally,
\[
 kZ_k=\sum_{|S|=k-1}\left(\prod_{i\in S}\lambda_i\right)
                           \sum_{j\notin S}\lambda_j.
\]
Removing any \(k-1\) indices removes at most the sum of the
largest \(k-1\) squared singular values.  Hence every inner
sum is at least \(T_k\), and
\(kZ_k\geq T_ke_{k-1}(\lambda)\).
This proves \eqref{four:volume-tail-operator}.
\end{proof}

\begin{lemma}[Projection control of volume marginals]
\label{four:volume-stability}
Let \(P\) be any orthogonal projection on the row-index space,
\(\ell_i=P_{ii}\), and
\(\varepsilon=\|(I-P)M\|_{\rm op}\).  With the notation of
Proposition~\ref{four:volume-operator},
\begin{equation}\label{four:volume-stability-tail}
 \sqrt{\alpha_i}
 \leq\sqrt{\ell_i}
      +\sqrt{\frac{k}{T_k}}\,\varepsilon\sqrt{1-\ell_i}.
\end{equation}
In particular for \(k=3\), \(\kappa=s_3(M)>0\),
\begin{equation}\label{four:volume-stability-three}
 \sqrt{\alpha_i}
 \leq\sqrt{\ell_i}
      +\frac{\sqrt3\,\varepsilon}{\kappa}\sqrt{1-\ell_i}.
\end{equation}
\end{lemma}
\begin{proof}
Use the triangle inequality after applying \(W_k^{1/2}\) to
\(e_i=Pe_i+(I-P)e_i\).  The bound \(W_k\leq I\) controls
the first component by \(\sqrt{\ell_i}\).
The first inequality in \eqref{four:volume-tail-operator}
controls the second by
\[
 \sqrt{k/T_k}\,\|M^{\mathsf T}(I-P)e_i\|
 \leq\sqrt{k/T_k}\,\varepsilon\sqrt{1-\ell_i}.
\]
Since \(\|W_k^{1/2}e_i\|=\sqrt{\alpha_i}\), this proves the
first assertion.  The bound \(T_3\geq\kappa^2\) proves the second.
\end{proof}

\begin{theorem}[A cross block near a diffuse range]
\label{four:near-diffuse}
Suppose a physical cross block of a quadratic form satisfies
\[
 M=M_0+E,\qquad s_3(M_0)\geq a>0,\qquad
 \max_i(P_{\operatorname{range}(M_0)})_{ii}\leq\tfrac12,
 \qquad \|E\|_{\rm op}\leq a/16.
\]
Then, with arbitrary self blocks, outside couplings, and affine
terms, the quadratic form has a continuous density vanishing
at infinity and
\begin{equation}\label{four:near-diffuse-bound}
 \|p\|_\infty\leq\frac{24\sqrt3\,e}{5a}<\frac{23}{a}.
\end{equation}
No rank bound on either \(M_0\) or \(E\) is required.
\end{theorem}
\begin{proof}
Let \(P=P_{\operatorname{range}(M_0)}\).
Then \(\varepsilon=\|(I-P)M\|\leq a/16\), and singular-value
perturbation gives \(\kappa=s_3(M)\geq15a/16\).
Lemma~\ref{four:volume-stability} gives
\[
 \sqrt{\alpha_i}\leq\sqrt{\ell_i}
                    +\frac{\sqrt3}{15}\sqrt{1-\ell_i}.
\]
For \(0\leq\ell\leq1/2\) and \(0\leq d\leq1\), the function
\(\sqrt\ell+d\sqrt{1-\ell}\) is increasing in \(\ell\).
Consequently
\[
 \alpha_i\leq\frac{(1+\sqrt3/15)^2}{2}<\frac23.
\]
For a rational check, replacing \(\sqrt3\) by \(2\) gives
the still smaller-than-\(2/3\) bound \(289/450\).
Theorem~\ref{four:volume-density} with \(k=3\), \(q=4/3\)
therefore gives
\[
 \|p\|_\infty\leq\frac{9\sqrt3\,e}{2}\,Z_3^{-1/6}.
\]
Since \(Z_3\geq s_1(M)^2s_2(M)^2s_3(M)^2\geq\kappa^6\),
the claimed constant follows from \(\kappa\geq15a/16\).
\end{proof}

\begin{corollary}[Distributed spectral tails]\label{four:tail}
Let \(M\) be a physical cross block of rank at least three and put
\[
 T_3=\sum_{j\geq3}s_j(M)^2.
\]
Suppose a projection \(P\) on its row-index space satisfies
\[
 \max_iP_{ii}\leq\tfrac12,\qquad
 \|(I-P)M\|_{\rm op}\leq\sqrt{T_3}/15.
\]
Then
\begin{align}
 \|p\|_\infty
 &\leq\frac{9\sqrt3\,e}{2}\,e_3(M^{\mathsf T}M)^{-1/6}
 \label{four:tail-bound}\\
 &\leq\frac{9\sqrt3\,e}{2}
           [s_1(M)^2s_2(M)^2T_3]^{-1/6}.
 \nonumber
\end{align}
If \(s_2(M)\geq a>0\) and \(T_3\geq b^2>0\), the latter
bound is \((9\sqrt3\,e/2)a^{-2/3}b^{-1/3}\).
In particular, when \(P\) projects onto the top two left singular
directions and its diagonal is at most \(1/2\), the condition
\(T_3/s_3(M)^2\geq225\) suffices.
\end{corollary}
\begin{proof}
Use \eqref{four:volume-stability-tail} with \(k=3\).
The same marginal estimate as in
Theorem~\ref{four:near-diffuse} gives \(\alpha_i<2/3\).
The fractional three-volume theorem proves the first bound.
The terms of \(e_3\) containing the first two indices sum to
\(s_1^2s_2^2T_3\), proving the second.  For the indicated
projection, \(\|(I-P)M\|_{\rm op}=s_3(M)\).
\end{proof}

\subsection{Small balls for an arbitrary nonzero quadratic form}

\begin{lemma}[Quarter-power concentration]\label{four:quarter}
Let \(S\ne0\) be any real symmetric matrix and let
\[
 R=X^{\mathsf T}SX+b^{\mathsf T}X+c.
\]
For every \(\rho>0\),
\begin{equation}\label{four:quarter-bound}
 \sup_w\mathbb P(|R-w|<\rho)
 \leq\min\left\{1,5\left(\frac{\rho}{\|S\|_{\rm op}}\right)^{1/4}\right\}.
\end{equation}
In particular \(R\) has no atoms.  The same concentration estimate
holds with a closed interval.
\end{lemma}
\begin{proof}
Put \(L=\|S\|_{\rm op}>0\) and \(m=\max_i|S_{ii}|\).
First suppose \(m\geq L/2\), and condition on all coordinates
except one attaining \(m\).
A real quadratic \(ax^2+dx+e\) with \(a\ne0\) has a preimage
of any interval of length \(2\rho\) of total length at most
\(2\sqrt{2\rho/|a|}\).  Indeed, complete the square and note
that for an interval \([u,u+\delta]\) the set where
\((x-x_0)^2\) belongs to that interval has length at most
\(2\sqrt\delta\): on the nonnegative half-line this is
\(2(\sqrt{u+\delta}-\sqrt u)\), with the evident truncation
when \(u<0\).
The density cap therefore bounds the conditional probability by
\(4\sqrt{\rho/L}\).
For \(\rho\leq L\), this is at most
\(4(\rho/L)^{1/4}\); for \(\rho>L\), the bound one suffices.

Now suppose \(m<L/2\).  Let \(S_0=S-\operatorname{diag}(S)\),
so \(\|S_0\|_{\rm op}\geq L/2\).
For independent uniform signs on the diagonal of a matrix \(E\),
define \(C_E=(S_0-ES_0E)/2\).
The zero diagonal of \(S_0\) implies
\(\mathbb EC_E=S_0/2\).  Convexity of the operator norm shows
that some sign choice satisfies
\(\|C_E\|_{\rm op}\geq\|S_0\|_{\rm op}/2\geq L/4\).
For the cut into positive and negative signs,
\[
 C_E=\begin{pmatrix}0&B\\B^{\mathsf T}&0\end{pmatrix},
 \qquad \|B\|_{\rm op}=\|C_E\|_{\rm op}\geq L/4 .
\]
This is a cut in the original coordinates.
The unrestricted determinant bound
\eqref{four:unrestricted-determinant}, retaining only the
largest singular value of \(B\), gives
\begin{equation}\label{four:quarter-decay}
 |\varphi_R(t)|
 \leq\left(1+\frac{4\|B\|_{\rm op}^2t^2}{\pi^2e^2}\right)^{-1/8}
 \leq\left(\frac{2\pi e}{L|t|}\right)^{1/4}.
\end{equation}
The affine terms are allowed throughout the decoupling argument.

For the concentration constant, let \(v>0\).  Then
\(\mathbf1_{\{|x|<\rho\}}\leq e^v e^{-vx^2/\rho^2}\).
Gaussian Fourier inversion and Fubini give
\[
 \mathbb P(|R-w|<\rho)
 \leq\frac{e^v\rho}{2\sqrt{\pi v}}
      \int_{\mathbb R}e^{-\rho^2t^2/(4v)}
                         |\varphi_R(t)|\,dt.
\]
For \(0<\alpha<1\),
\[
 \int_{\mathbb R}e^{-at^2}|t|^{-\alpha}\,dt
 =a^{(\alpha-1)/2}\Gamma((1-\alpha)/2),
\]
by the substitution \(u=at^2\).
Use \eqref{four:quarter-decay}, \(\alpha=1/4\), and \(v=1/8\).
The coefficient of \((\rho/L)^{1/4}\) is
\[
 2^{3/8}\pi^{-1/4}e^{3/8}\Gamma(3/8)<5.
\]
For an elementary verification,
\(\Gamma(3/8)\leq8/3+e^{-1}<73/24\), obtained by splitting
its defining integral at one.  Also
\(2^{3/8}<4/3\), \(\pi^{-1/4}<4/5\), and \(e^{3/8}<3/2\).
The product of these four upper bounds is \(73/15<5\).
This proves \eqref{four:quarter-bound}.
Letting \(\rho\downarrow0\) proves absence of atoms.
The closed-interval version follows by applying the result with
\(\rho+\varepsilon\) and then letting \(\varepsilon\downarrow0\).
\end{proof}

\begin{corollary}[Logarithmic and fractional potentials]
\label{four:log-potential}
In the setting of Lemma~\ref{four:quarter}, let
\(L=\|S\|_{\rm op}\).  For every \(a>0\), put
\(M=5(a/L)^{1/4}\).  Uniformly in \(w\),
\begin{align}
 \mathbb E\log_+\frac a{|R-w|}
 &\leq
 \begin{cases}
  4M,&M\leq1,\\
  4+4\log M,&M\geq1,
 \end{cases}
 \label{four:log-potential-bound}\\
 &\leq20(a/L)^{1/4}.
 \nonumber
\end{align}
For every \(0<\beta<1/4\), moreover,
\begin{equation}\label{four:negative-moment}
 \sup_w\mathbb E|R-w|^{-\beta}
 \leq\frac{5^{4\beta}}{1-4\beta}\,L^{-\beta}.
\end{equation}
\end{corollary}
\begin{proof}
Tonelli's theorem and the quarter-power bound give
\[
 \mathbb E\log_+\frac a{|R-w|}
 =\int_0^\infty
      \mathbb P(|R-w|<ae^{-u})\,du
 \leq\int_0^\infty\min\{1,Me^{-u/4}\}\,du.
\]
If \(M\leq1\), this integral is \(4M\).
If \(M>1\), split at \(u=4\log M\) to obtain \(4+4\log M\).
Both cases are at most \(4M\).
For the negative moment use
\[
 \mathbb E|R-w|^{-\beta}
 =\beta\int_0^\infty r^{-\beta-1}
                       \mathbb P(|R-w|<r)\,dr.
\]
Split at \(r_0=L/5^4\).  The probability is at most
\((r/r_0)^{1/4}\) for \(r<r_0\) and at most one otherwise.
The two integrals sum to
\(r_0^{-\beta}/(1-4\beta)\), giving
\eqref{four:negative-moment}.
\end{proof}

The spectral-window, volume, and ridge criteria above are
quantitative sufficient hypotheses.  A lower bound for the third
singular value of the full quadratic matrix alone does not verify
their coordinate cut or marginal conditions.  Their usefulness is
that the estimates retain arbitrary self interactions and, where
specified, all additional spectral directions without any assumption
that a remainder is independent of the principal part.
